%%%%%%%%%%%%%%%%%%%%%%%%%%%%%%%%%%%%%%%%%
% Plantilla de tesis FaMAF/UNC
% Archivo principal corregido para Overleaf
%%%%%%%%%%%%%%%%%%%%%%%%%%%%%%%%%%%%%%%%%

\documentclass[
  11pt,
  oneside,
  english,
  spanish,
  singlespacing,
  headsepline
]{MastersDoctoralThesis}

\usepackage[utf8]{inputenc}
\usepackage[T1]{fontenc}
\usepackage{mathpazo}
\usepackage{amsmath}
\IfFileExists{biblatex.sty}{%
  \usepackage[backend=bibtex,style=numeric,natbib=true,bibencoding=utf-8]{biblatex}
  \addbibresource{list.bib}
}{%
  % Respaldo para instalaciones mínimas. Overleaf normalmente usa biblatex.
  \usepackage[numbers]{natbib}
  \newcommand{\printbibliography}{%
    \bibliographystyle{plainnat}%
    \bibliography{list}%
  }
}
\usepackage[autostyle=true]{csquotes}

\geometry{
  paper=a4paper,
  inner=2.5cm,
  outer=3.8cm,
  bindingoffset=.5cm,
  top=1.5cm,
  bottom=1.5cm
}
\setlength{\headheight}{13.6pt}

\thesistitle{Restauración de imágenes mediante MAXIM y Mixture of Experts}
\author{Matías A. Murat}
\supervisor{Martín A. Domínguez y Francisco A. Tamarit}
\subject{}
\keywords{}
\university{Universidad Nacional de Córdoba}
\department{}
\faculty{Facultad de Matemática, Astronomía, Física y Computación}

\AtBeginDocument{
  \hypersetup{
    pdftitle={\ttitle},
    pdfauthor={\authorname},
    pdfkeywords={\keywordnames}
  }
}

\begin{document}
\sloppy
\frontmatter
\pagestyle{plain}

\hypersetup{pageanchor=false}
\begin{titlepage}
\begin{center}
  \IfFileExists{Figures/crest.png}{%
    \includegraphics[scale=0.5]{Figures/crest.png}\par
    \vspace*{.06\textheight}
  }{}

  {\scshape\LARGE \facname\par}\vspace{1cm}
  \textsc{\Large Trabajo especial de Licenciatura en Ciencias de la Computación}\\[0.5cm]

  \HRule\\[0.4cm]
  {\huge\bfseries \ttitle\par}\vspace{0.4cm}
  \HRule\\[1.5cm]

  \begin{minipage}[t]{0.4\textwidth}
    \begin{flushleft}\large
      \emph{Autor:}\\
      \authorname
    \end{flushleft}
  \end{minipage}
  \begin{minipage}[t]{0.4\textwidth}
    \begin{flushright}\large
      \emph{Directores:}\\
      \supname
    \end{flushright}
  \end{minipage}\\[3cm]

  \vfill
  {\large \today}\\[2cm]

  \IfFileExists{licenciagv.png}{\includegraphics[scale=0.8]{licenciagv.png}\par}{}
  Esta obra está bajo una
  \href{https://creativecommons.org/licenses/by-nc-sa/4.0/}%
       {Licencia Creative Commons Atribución-NoComercial-CompartirIgual 4.0 Internacional}.
  \vfill
\end{center}
\end{titlepage}
\hypersetup{pageanchor=true}

\begin{abstract}
\addchaptertocentry{Resumen}
Acá va el resumen. Debe ser breve, aproximadamente media carilla.
\end{abstract}

\begin{acknowledgements}
\addchaptertocentry{Agradecimientos}
Acá van los agradecimientos. Esta sección es optativa.
\end{acknowledgements}

{
  \hypersetup{linkcolor=black}
  \tableofcontents
  \listoffigures
  \renewcommand{\listtablename}{Índice de tablas}
  \listoftables
}

\dedicatory{Dedicado a quien quieras. Esta página es optativa.}

\mainmatter
\pagestyle{thesis}

\chapter{Introducción}
\label{chap:introduccion}

La producción y el intercambio de imágenes digitales forman parte de una gran
cantidad de actividades cotidianas, científicas e industriales. Las imágenes
capturadas por cámaras, teléfonos, sistemas de vigilancia, instrumental médico
o sensores remotos rara vez se obtienen en condiciones ideales. El movimiento
de la cámara o de la escena puede introducir desenfoque; la iluminación
insuficiente reduce el contraste y oculta detalles; las condiciones
atmosféricas alteran la visibilidad; el ruido del sensor degrada las regiones de
baja señal y la lluvia puede superponer estructuras ajenas a la escena. Estas
alteraciones no solo afectan la calidad visual, sino que también pueden
perjudicar tareas posteriores de visión por computadora, como la detección, la
segmentación o el reconocimiento de objetos.

La restauración de imágenes estudia cómo recuperar una imagen limpia a partir
de una observación degradada. El problema es difícil porque diferentes imágenes
limpias pueden dar lugar a degradaciones similares y porque, en aplicaciones reales,
el tipo y la intensidad de la degradación no siempre son conocidos. Por esta razón,
una reconstrucción adecuada requiere aprovechar tanto la información local: bordes, texturas y
detalles finos, así como el contexto global de la imagen.

Durante los últimos años, los métodos basados en aprendizaje profundo han
reemplazado gran parte de los procedimientos diseñados de manera específica
para cada degradación. A partir de pares de imágenes degradadas y sus
correspondientes referencias, una red neuronal puede aprender una transformación
de restauración, sin embargo, es habitual entrenar un modelo distinto para
cada tarea. Esta estrategia permite una alta especialización, pero obliga a
mantener y ejecutar varios modelos cuando un sistema debe responder correctamente
ante degradaciones variadas.

Una alternativa consiste en entrenar un único modelo para varias tareas de
restauración. Este enfoque reduce la duplicación de parámetros y permite
compartir representaciones entre problemas relacionados. Al mismo tiempo,
introduce una dificultad: eliminar ruido, compensar desenfoque, remover lluvia,
reducir neblina y realzar una imagen oscura requieren transformaciones
diferentes y, en algunos casos, parcialmente opuestas. Un único camino de
cómputo debe aprender entonces una solución suficientemente general para todas
ellas. Si la capacidad compartida no puede adaptarse al tipo de degradación
presente en cada entrada, el aprendizaje conjunto puede producir interferencia
entre tareas y limitar la calidad de la reconstrucción.

En este trabajo se adopta MAXIM (\textit{Multi-Axis MLP for Image Processing}) como 
arquitectura base para la restauración de imágenes \citep{tu2022maxim}. Su diseño sigue una estructura codificador/decodificador tipo U-Net, con conexiones \textit{skip} que combinan características de bajo y alto nivel y contribuyen a preservar la información espacial. A diferencia de las arquitecturas puramente convolucionales, MAXIM incorpora bloques \textit{Multi-Axis Gated MLP}, que capturan interacciones locales  y globales. Esta combinación permite integrar información en distintas escalas y capturar tanto detalles finos como patrones globales. Sin embargo, en su formulación estándar, todas las imágenes atraviesan el mismo flujo de procesamiento, independientemente del tipo de degradación presente, lo que motiva estudiar mecanismos de especialización para un escenario de restauración multi-tarea.

La hipótesis que motiva esta tesis es que dicho flujo compartido puede
complementarse con el paradigma \textit{Mixture of Experts} (MoE). En una mezcla
de expertos, varios módulos con capacidad de especialización colaboran bajo el
control de un ruteador, que asigna a cada experto un peso en base a 
la entrada \citep{jacobs1991adaptive,shazeer2017outrageously}. En lugar de exigir
que todos los parámetros resuelvan de la misma forma cualquier degradación, el
modelo puede compartir una parte sustancial del procesamiento y reservar otra
parte para respuestas especializadas. El ruteador actúa como un mecanismo de
adaptación: a partir de las características de la imagen, determina qué expertos
deben intervenir y con qué peso en la reconstrucción.

La combinación de MAXIM y MoE plantea un compromiso entre generalización,
especialización y costo computacional. Replicar una red completa por cada tarea
produciría especialistas independientes, pero eliminaría gran parte del
beneficio de un modelo unificado. En el extremo opuesto, reemplazar por completo
la salida de MAXIM con expertos inicializados de manera similar puede dificultar
la diferenciación entre ellos y perder una referencia común de reconstrucción.
La arquitectura estudiada en este trabajo conserva, por lo tanto, una cabeza
compartida equivalente a la salida original de MAXIM y utiliza los expertos para
modelar correcciones residuales condicionadas por la entrada. Si estas
correcciones son nulas, el modelo recupera exactamente el camino compartido de
la arquitectura base.

La ubicación de los expertos, su capacidad, la señal utilizada por el ruteador
y la forma de entrenar la mezcla continúan siendo decisiones centrales. En
particular, el ruteador puede activar todos los expertos o seleccionar solamente
un subconjunto mediante un esquema \textit{top-k}. También es necesario
preservar los mecanismos de supervisión profunda y multi-escala de MAXIM. En la
implementación considerada, las salidas auxiliares permanecen conectadas a sus
pérdidas originales y la mezcla de expertos modifica únicamente la predicción
final a resolución completa.

El problema abordado puede formularse del siguiente modo: dada una colección
de pares $(x,y)$, donde $x$ es una imagen afectada por alguna de varias
degradaciones e $y$ es su referencia limpia, se busca aprender una función
$F$ tal que $F(x)$ aproxime a $y$ sin requerir, como condición de entrada, una
etiqueta que indique la tarea. Para la variante con expertos, esta función se
expresa conceptualmente como

\begin{equation}
  \widehat{y} = x {}+ H_{\mathrm{comp}}\!\left(f_{\theta}(x)\right)
  {}+ \sum_{k=1}^{K} p_k\!\left(f_{\theta}(x)\right)
      \,\Delta E_k\!\left(f_{\theta}(x)\right),
  \label{eq:intro-moe}
\end{equation}

donde $f_{\theta}$ representa las características compartidas producidas por
MAXIM, $H_{\mathrm{comp}}$ es la cabeza de reconstrucción compartida y
$\Delta E_k$ es la corrección propuesta por el $k$-ésimo experto. El término
$x$ conserva la conexión residual entre la imagen degradada y la predicción
final. Los pesos $p_k$ representan las asignaciones finales del ruteador,
calculadas a partir de las características de la entrada, y satisfacen
$p_k\geq 0$ y $\sum_{k=1}^{K}p_k=1$. En un ruteo \textit{top-k}, como máximo
$k$ de estos pesos son distintos de cero. La Ecuación~\ref{eq:intro-moe}
describe la idea general; la
capacidad de las cabezas, la forma concreta del ruteador y las pérdidas
empleadas se presentan en los capítulos metodológicos.

El objetivo general de esta tesis es diseñar y evaluar una extensión de MAXIM
basada en mezcla de expertos para restaurar imágenes afectadas por distintos
tipos de degradación dentro de un único modelo. La evaluación busca determinar
si la especialización aporta ventajas frente a una versión base MAXIM multi-tarea entrenada en condiciones equivalentes.

Para alcanzar este objetivo general se plantean los siguientes objetivos
específicos:

\begin{itemize}
  \item implementar y entrenar una versión base de MAXIM que sirva como referencia
        reproducible para las comparaciones;
  \item construir un conjunto de entrenamiento multi-tarea con ejemplos de
        eliminación de desenfoque, neblina, ruido y lluvia, además de realce de
        imágenes con baja iluminación;
  \item incorporar una cabeza compartida y expertos residuales a la arquitectura,
        conservando la supervisión profunda y el procesamiento en múltiples
        escalas;
  \item estudiar estrategias de ruteo y regularización que favorezcan el uso
        diferenciado de los expertos sin provocar una asignación degenerada;
  \item comparar las variantes mediante métricas cuantitativas de calidad,
        inspección visual, costo computacional y análisis del comportamiento del
        ruteador.
\end{itemize}

La comparación experimental se diseña para aislar, en la medida de lo posible,
el efecto de la mezcla de expertos. Esto exige utilizar los mismos datos, el
mismo presupuesto de entrenamiento y configuraciones de capacidad comparables
entre el modelo propuesto y la línea base. Además de medir la calidad final de
las imágenes, es necesario observar si los expertos desarrollan funciones
diferenciadas y si las decisiones del ruteador guardan relación con las
degradaciones presentes. De este modo, el análisis no queda limitado a una
única métrica agregada, sino que considera también el mecanismo que se pretende
introducir.

El resto de la tesis se organiza de la siguiente manera. En primer lugar se
presentan los fundamentos de restauración de imágenes, MAXIM y las mezclas de
expertos. A continuación se describen las herramientas utilizadas, los conjuntos
de datos, la arquitectura propuesta, la implementación y el protocolo
experimental. Finalmente, se exponen y discuten los resultados de las
comparaciones, para luego presentar las conclusiones y posibles líneas de
continuación. Esta organización mantiene separados el planteo del problema, las
decisiones metodológicas y la evidencia experimental utilizada para evaluar la
propuesta.

\chapter{Fundamentos teóricos}
\label{chap:fundamentos-teoricos}

\chapter{Herramientas}
\label{chap:herramientas}

En este capítulo se describen brevemente las tecnologías empleadas para
implementar los modelos, construir el flujo de datos, ejecutar los
entrenamientos y analizar los resultados. La separación entre herramientas de
cómputo numérico, definición de redes neuronales, carga de datos y entorno de
ejecución resulta especialmente relevante en este trabajo: aunque todas
participan del mismo experimento, cada una cumple una función distinta dentro
del sistema.

\section{Lenguaje y cómputo numérico}

\subsection{Python}

La implementación se desarrolló en Python, un lenguaje de propósito general con
un ecosistema amplio para cómputo científico y aprendizaje automático
\citep{pythonDocs}. Su sintaxis y su integración con entornos interactivos
permitieron organizar en un mismo flujo la preparación de los datos, la
definición de los modelos, el entrenamiento, la evaluación y la generación de
diagnósticos.

El código se distribuyó principalmente en módulos reutilizables y cuadernos de
experimentación. Los módulos contienen la arquitectura MAXIM, las cabezas
expertas, el ruteador y las funciones comunes, mientras que los cuadernos
configuran cada corrida, montan los datos, ejecutan validaciones previas y
registran métricas y checkpoints.

\subsection{NumPy}

NumPy proporciona arreglos multidimensionales y operaciones vectorizadas para
cómputo científico en Python \citep{harris2020numpy}. En este trabajo se utilizó
como representación intermedia para resultados transferidos desde el pipeline
de datos y para operaciones auxiliares de evaluación. Entre ellas se encuentran
la acumulación de métricas, el cálculo de promedios globales y por tarea, y la
preparación de datos para visualizaciones.

\section{Aprendizaje profundo}

\subsection{JAX}

JAX es una biblioteca de cómputo numérico orientada a transformaciones de
programas sobre arreglos \citep{jax2018}. Ofrece diferenciación automática,
vectorización y compilación justo a tiempo (JIT) mediante \texttt{jax.jit}. Esta
última permite compilar las funciones de entrenamiento y ejecutarlas sobre
aceleradores, reduciendo el costo de las operaciones repetidas durante las
épocas.

La implementación de MAXIM utilizada como base fue desarrollada en este
ecosistema. JAX se empleó para las operaciones tensoriales del modelo, el
cálculo de gradientes, la actualización durante el entrenamiento y la
evaluación de las reconstrucciones.

\subsection{Flax}

Flax es una biblioteca de redes neuronales construida sobre JAX
\citep{flax2023}. La implementación utiliza su API \textit{Linen}, en la que las
capas y los modelos se expresan como subclases de \texttt{nn.Module}. Este
esquema se usó tanto para los bloques originales de MAXIM como para las
extensiones de esta tesis: la cabeza compartida, los expertos residuales y el
ruteador.

Flax también organiza los parámetros y el estado del modelo en estructuras
compatibles con las transformaciones de JAX. Esto permitió conservar el camino
de reconstrucción de MAXIM, incorporar módulos especializados y restaurar
checkpoints sin acoplar la definición de la arquitectura a un cuaderno
particular.

En conjunto, la implementación se realizó en Python utilizando JAX y el framework Flax,
aprovechando la compilación JIT y la diferenciación automática. Para acotar el consumo de memoria
en la variante MoE se recurrió al \textit{gradient checkpointing}, es decir, a la 
recomputación de activaciones durante el paso hacia atrás, en vez de almacenarlas en memoria.
Los entrenamientos se ejecutaron en una única GPU de Colab con $22\,\mathrm{GB}$ de VRAM; esta
restricción de memoria motivó tanto el tamaño del \textit{batch} como el diseño
de expertos de salida, en lugar de replicar bloques internos del
\textit{backbone}.

\section{Carga y procesamiento de datos}

\subsection{TensorFlow Data}

Aunque los modelos se implementaron con JAX y Flax, los datos de entrada se cargaron y procesaron con 
TensorFlow \citep{tensorflow2016} utilizando la API \texttt{tf.data} del framework.
Esta API permite representar colecciones potencialmente grandes como secuencias
de transformaciones, aplicar operaciones en paralelo, agrupar ejemplos en
batches y anticipar su carga mediante \textit{prefetch}.

Para cada tarea de restauración se construyó un conjunto de pares, cada uno
formado por una imagen degradada y su correspondiente imagen de referencia.
Durante el entrenamiento, cada conjunto de datos se repite de manera
independiente antes de combinarse con los demás mediante
\texttt{sample\_from\_datasets}. De este modo, las cinco tareas permanecen
disponibles durante toda la época y sus ejemplos pueden muestrearse uniformemente, 
independientemente del tamaño de cada dataset. El mismo pipeline aplica recortes y 
aumentos alineados a ambos elementos del par. La validación, en cambio, concatena los
conjuntos sin muestreo y recorre de manera determinista todos los ejemplos.

\subsection{Pillow}

Pillow es una biblioteca de procesamiento de imágenes para Python con soporte
para distintos formatos de archivos \citep{pillowDocs}. Se utilizó para abrir y
decodificar las imágenes de los conjuntos de datos, convertirlas a tres canales
de color (RGB) y realizar pruebas de integridad. También formó parte de las
pruebas que verifican que cada ruta listada pueda convertirse en un array con
las dimensiones y el rango esperados antes de iniciar un entrenamiento largo.

\section{Ejecución y almacenamiento}

\subsection{Google Colaboratory}

Google Colaboratory es un servicio alojado de cuadernos Jupyter que proporciona
acceso a recursos de cómputo, incluidas GPU y TPU, sin requerir una instalación
local completa \citep{googleColab}. Los entrenamientos y las evaluaciones de
esta tesis se ejecutaron en sesiones de Colab con GPU. Los cuadernos permitieron
reunir la configuración de cada experimento, las comprobaciones previas, la
ejecución y una primera inspección de los resultados.

Los recursos asignados por Colab pueden variar entre sesiones y el entorno de
ejecución es temporal. Por este motivo, las versiones de JAX y Flax se fijaron
explícitamente y cada corrida se guardó en un directorio independiente. También
se conservaron el archivo de configuración, el historial de entrenamiento, los
diagnósticos del ruteador y los checkpoints necesarios para reanudar o evaluar
un modelo.

\subsection{Google Drive}

Google Drive se utilizó como almacenamiento persistente vinculado a los
cuadernos de Colab. Allí se mantuvieron los datasets, las listas de
entrenamiento y validación, los checkpoints y los resultados de cada
experimento. Para reducir el riesgo de utilizar una versión incompleta de los
datos, las listas principales se identificaron mediante conteos y resúmenes
SHA-256, y los cuadernos ejecutaron pruebas de lectura sobre batches reales
antes de habilitar el entrenamiento.

\section{Visualización de resultados}

\subsection{Matplotlib}

Matplotlib es una biblioteca de visualización para Python
\citep{hunter2007matplotlib}. Se empleó para representar la evolución del loss y 
del PSNR, comparar el desempeño por tarea y construir mapas de
calor del uso de los expertos. Estas visualizaciones complementan las métricas
agregadas: permiten detectar diferencias entre tareas y observar si el ruteador
concentra sus decisiones, distribuye la carga o utiliza de manera efectiva los
expertos disponibles.

\appendix
\printbibliography
\end{document}
